\documentclass[a4paper,11pt]{article}

% --- Pacchetti Fondamentali ---
\usepackage[utf8]{inputenc}
\usepackage[T1]{fontenc}
\usepackage[italian]{babel}
\usepackage{geometry}
\geometry{top=2.5cm, bottom=2.5cm, left=2.5cm, right=2.5cm}

% --- Font Moderni ---
\usepackage{lmodern}
\usepackage{helvet}
\renewcommand{\familydefault}{\sfdefault} 

% --- Colori e Grafica ---
\usepackage{xcolor}
\usepackage{tikz}
\usetikzlibrary{shapes, arrows.meta, positioning, chains, fit, calc, matrix, trees, backgrounds}

% --- Definizione Colori (Light Theme + Orange + New Cover Colors)
\definecolor{mainOrange}{RGB}{255, 140, 0}
\definecolor{jsYellow}{HTML}{F7DF1E}
\definecolor{codeBg}{gray}{0.95}
\definecolor{codeBlue}{HTML}{0000FF}
\definecolor{codePink}{HTML}{800080}
\definecolor{codeGreen}{HTML}{008000}
\definecolor{niceBlue}{HTML}{2E86C1}
\definecolor{jsDark}{HTML}{282C34}    
\definecolor{tsBlue}{HTML}{3178C6}    
\definecolor{textGray}{RGB}{100, 100, 100}

% --- PALETTE COLORI ACCADEMICA (Nuovo Capitolo) ---
\definecolor{paletteIndigo}{HTML}{4B0082}
\definecolor{paletteRegalia}{HTML}{522D80}
\definecolor{paletteLenurple}{HTML}{BA55D3}
\definecolor{paletteLavender}{HTML}{E6E6FA}

% --- COMANDO COPERTINA ESERCIZIARIO (Pure TikZ + Optical Centering) ---
\newcommand{\jscover}[1]{
    \begin{titlepage}
    \thispagestyle{empty}
    \begin{tikzpicture}[remember picture, overlay]
        
        % 1. Sfondo Laterale (Scuro)
        \fill[jsDark] (current page.north west) rectangle ($(current page.south west) + (4cm,0)$);
        
        % 2. Linea Accento (Arancione)
        \draw[mainOrange, line width=4pt] ($(current page.north west) + (4.2cm,-2cm)$) -- ($(current page.south west) + (4.2cm,2cm)$);
        
        % 3. Contenuto Testuale (TikZ Node Absolute Positioning)
        % Optical Centering: Width 12cm positioned at x=12.5cm
        \node[anchor=north, inner sep=0, text width=12cm, align=center] at ($(current page.north west) + (12.5cm,-5cm)$) {
            
                
                {\fontsize{18}{22}\selectfont \textcolor{textGray}{Università di Pisa}}\\[0.2cm]
                {\fontsize{16}{20}\selectfont \textcolor{textGray}{Corso di Laurea in Informatica}}\\[2cm]
                
                % RIDOTTO FONT PER FARLO STARE IN 12CM
                {\fontsize{42}{52}\selectfont \textbf{\textcolor{jsDark}{L'ESERCIZIARIO}}}\\[0.5cm]
                
                {\fontsize{35}{45}\selectfont \textbf{\textcolor{black}{DI LABORATORIO}}}\\[0.2cm]
                {\fontsize{35}{45}\selectfont \textbf{\textcolor{black}{JAVASCRIPT}}}\\[1.5cm]
                
                \textcolor{mainOrange}{\rule{10cm}{4pt}}\\[1cm]
                
                {\Large \textbf{Dal Testo alla Logica}}\\[0.2cm]
                {\large \textit{Oltre 48 Esercizi Svolti e Commentati}}\\[1.5cm]
                
                {\large Docente titolare:}\\[0.2cm]
                {\Large \textbf{#1}}\\[2.5cm]
                
                \textbf{\large Autore del riassunto:}\\
                {\large Joseph Zucchelli (IlBroccolo)}\\[2cm]
                
                \vfill
                {\large \textbf{Anno Accademico 2025 - 2026}}
            
        };
        
        % Decorazione JS + TS
        \node[anchor=south east] at ($(current page.south east) + (-2cm, 2cm)$) {
            \tikz\draw[fill=jsYellow, draw=none] (0,0) rectangle (2,2) node[midway, black, font=\bfseries\huge] {JS};
             \tikz\draw[fill=tsBlue, draw=none] (2.2,0) rectangle (4.2,2) node[midway, white, font=\bfseries\huge] {TS};
        };

    \end{tikzpicture}
    \end{titlepage}
    \clearpage
}

% --- Box per il Codice e Note (tcolorbox) ---
\usepackage[most]{tcolorbox}
\usepackage{listings}
\usepackage{titlesec}

% Setup Titoli
\titleformat{\section}
{\color{mainOrange}\normalfont\Large\bfseries}
{\color{mainOrange}\thesection}{1em}{}[\color{mainOrange}\hrule height 1pt]

\titleformat{\subsection}
{\color{mainOrange}\normalfont\large\bfseries}
{\color{mainOrange}\thesubsection}{1em}{}

% Stile per il codice JS/TS
\lstdefinelanguage{JavaScript}{
  keywords={typeof, new, true, false, catch, function, return, null, catch, switch, var, if, in, while, do, else, case, break, const, let, interface, type, number, string, boolean, any, unknown, void},
  keywordstyle=\color{blue}\bfseries,
  ndkeywords={class, export, boolean, throw, implements, import, this, console, log},
  ndkeywordstyle=\color{teal}\bfseries,
  identifierstyle=\color{black}, 
  sensitive=false,
  comment=[l]{//},
  morecomment=[s]{/*}{*/},
  commentstyle=\color{green!50!black}\ttfamily,
  stringstyle=\color{red!60!black}\ttfamily,
  morestring=[b]',
  morestring=[b]"
}

% Box Codice (JSBox)
\newtcblisting{jsbox}[1][]{
  colback=codeBg,
  colframe=gray!50,
  listing only,
  listing options={
    language=JavaScript,
    basicstyle=\ttfamily\small,
    numbers=left,
    numberstyle=\tiny\color{gray},
    breaklines=true,
    columns=fullflexible
  },
  title={\textbf{Soluzione / Codice}},
  coltitle=black,
  fonttitle=\bfseries\ttfamily,
  enhanced,
  attach boxed title to top right={yshift=-3mm, xshift=-3mm},
  boxed title style={colback=jsYellow},
  boxrule=1pt,
  #1
}

% Box Concetti (Dictionary)
\newtcolorbox{dizionario}[1]{
  colback=niceBlue!5,
  colframe=niceBlue,
  title={#1},
  fonttitle=\bfseries,
  sharp corners=downhill,
  enhanced,
  drop shadow
}

% Box Esercizio
\newtcolorbox{esercizio}[1]{
  colback=white,
  colframe=mainOrange,
  title={#1},
  fonttitle=\bfseries,
  enhanced,
  attach boxed title to top left={yshift=-3mm, xshift=3mm},
  boxed title style={colback=mainOrange},
  boxrule=1pt
}

% --- NUOVI STILI BOX (Per capitoli esterni) ---
\newtcolorbox{definition}[1][]{
    colback=paletteLavender!30,
    colframe=paletteIndigo,
    fonttitle=\bfseries,
    title={Definizione: #1},
    enhanced,
    breakable,
    attach boxed title to top left={yshift=-2mm, xshift=2mm},
    boxed title style={colback=paletteIndigo}
}

\newtcolorbox{concept}[1][]{
    colback=white,
    colframe=paletteRegalia,
    fonttitle=\bfseries,
    title={Concetto Chiave: #1},
    enhanced,
    breakable,
    borderline west={2pt}{0pt}{paletteRegalia}
}

\newtcolorbox{curiosita}[1][]{
    colback=paletteLenurple!10,
    colframe=paletteLenurple,
    fonttitle=\bfseries,
    title={Nota a margine: #1},
    enhanced,
    breakable
}

\newtcolorbox{warning}[1][]{
    colback=red!5,
    colframe=red!60!black,
    fonttitle=\bfseries,
    title={Attenzione: #1},
    enhanced,
    breakable
}

\begin{document}

\jscover{Prof.ssa Ferrari}
\tableofcontents
\newpage

% --- Capitoli ---
\section{Capitolo 0: Dal Testo alla Logica}

In questo capitolo preliminare, impariamo a decifrare il "linguaggio da esame". Spesso la difficoltà non è scrivere il codice, ma capire \textit{esattamente} cosa viene chiesto.
Usiamo il \textbf{Dizionario Tecnico} per tradurre le parole chiave in codice.

\subsection{1. Il Glossario delle Strutture Dati}

\begin{dizionario}{Array Annidato (Nested)}
\textbf{Definizione:} Un array che contiene al suo interno altri array.
\textbf{Esempio:} \texttt{[1, [2, 3], 4]}.
\textbf{Cosa Significa nel Codice:} Probabilmente dovrai usare la \textbf{ricorsione} (vedi `flat` o `partition`) o doppi cicli.
\end{dizionario}

\begin{dizionario}{Insieme (Set) implementato come Oggetto}
\textbf{Definizione:} Un oggetto dove la chiave è l'elemento e il valore è \texttt{true}.
\textbf{Codice:} \texttt{\{ "mela": true, "pera": true \}}.
\textbf{Perché?} Per avere un accesso veloce $O(1)$ invece di cercare nell'array $O(n)$.
\end{dizionario}

\begin{dizionario}{Multi-insieme}
\textbf{Definizione:} Come un Set, ma contiamo quante volte appare un elemento.
\textbf{Codice:} \texttt{\{ "mela": 2, "pera": 1 \}}.
\textbf{Pattern Tipico:} \texttt{if (k in obj) obj[k]++; else obj[k] = 1;}
\end{dizionario}

\subsection{2. Il Glossario per le Funzioni}

\begin{dizionario}{Predicato}
\textbf{Definizione:} Una funzione che restituisce \textbf{sempre e solo} \texttt{boolean}.
\textbf{Firma:} \texttt{f: (x: T) => boolean}
\textbf{Esempio:} \texttt{isPari(x)}, \texttt{isMaggiorenne(user)}.
\end{dizionario}

\begin{dizionario}{Funzione Fabbrica (Factory)}
\textbf{Definizione:} Una funzione che \textbf{crea e restituisce} una nuova funzione.
\textbf{Concetto Chiave:} \textbf{Closure}. La funzione restituita "ricorda" i parametri della fabbrica.
\end{dizionario}

\begin{dizionario}{In-Place (Side Effect)}
\textbf{Definizione:} Modificare direttamente l'oggetto/array passato come argomento, SENZA crearne uno nuovo.
\textbf{Attenzione:} Se il testo dice "senza side-effect" o "funzionale pura", devi clonare l'input prima di toccarlo!
\end{dizionario}

\subsection{Esercizi di "Traduzione"}

\begin{esercizio}{Traduzione 1: La "Fabbrica di Predicati"}
\textbf{TESTO:} "Scrivere una funzione `createFilter` che accetta un numero `soglia` e restituisce un predicato che vale true se l'input è maggiore della soglia."

\textbf{ANALISI:}
1. \textit{Scrivere una funzione `createFilter` che accetta `soglia`}: \texttt{function createFilter(soglia) \{ ... \}}
2. \textit{Restituisce un predicato}: Deve fare \texttt{return function(x) \{ ... \}}
3. \textit{Vale true se input > soglia}: \texttt{return x > soglia}

\end{esercizio}

\begin{jsbox}
function createFilter(soglia) {
    // Ritorna una funzione (Closure) che ricorda 'soglia'
    return function(valoreInput) {
        return valoreInput > soglia; // Ritorna boolean (Predicato)
    };
}

const filtroMaggiorenni = createFilter(18);
console.log(filtroMaggiorenni(20)); // true
\end{jsbox}

\begin{esercizio}{Traduzione 2: "Senza Funzioni di Libreria"}
\textbf{TESTO:} "Calcolare l'intersezione tra due array A e B senza usare `Set`, `filter` o `includes`."

\textbf{ANALISI:}
1. \textit{Intersezione}: Elementi comuni.
2. \textit{Senza `includes`}: Devi farti un doppio ciclo for a mano o usare un oggetto di appoggio (Dictionary) per la velocità.
\end{esercizio}

\begin{jsbox}
function intersezioneManuale(A, B) {
    let result = [];
    // Strategia Dizionario per efficienza
    let mappaA = {};
    
    // 1. Popolo il dizionario con gli elementi di A
    for (let i = 0; i < A.length; i++) {
        mappaA[A[i]] = true;
    }
    
    // 2. Controllo se gli elementi di B sono nel dizionario
    for (let j = 0; j < B.length; j++) {
        let val = B[j];
        if (mappaA[val]) { // Equivale a A.includes(val)
            result.push(val);
        }
    }
    return result;
}
\end{jsbox}

\subsection{Tabella Rapida Concetti Matematici}
\begin{enumerate}
    \item \textbf{Suriettiva:} Copre tutto il codominio.
    \item \textbf{Iniettiva:} Nessuna collisione (input diversi $\to$ output diversi).
    \item \textbf{Invertibile:} Biunivoca (1-a-1 perfetto).
    \item \textbf{Punto Fisso:} Quando $f(x) = x$.
\end{enumerate}
 \newpage
\section{Capitolo 1: Strutture Dati (Array, Set, Alberi)}

\begin{esercizio}{Es. 1: Il "Flatten" Manuale (Array Annidati)}
\textbf{TRACCIA:} Scrivere una funzione ricorsiva `appiattisci(arr)` che prende un array che può contenere numeri o altri array (annidamento infinito) e restituisce un unico array piatto con tutti i numeri.
\textit{Non usare `flat()` nativo.}

\textbf{Esempio:} `[1, [2, [3, 4]], 5]` $\to$ `[1, 2, 3, 4, 5]`
\end{esercizio}

\begin{jsbox}
function appiattisci(arr) {
    let risultato = [];
    
    for (let i = 0; i < arr.length; i++) {
        let el = arr[i];
        
        if (Array.isArray(el)) {
            // Se è un array, chiamata ricorsiva!
            let sottoArrayPiatto = appiattisci(el);
            // Concateno i risultati
            for (let j = 0; j < sottoArrayPiatto.length; j++) {
                risultato.push(sottoArrayPiatto[j]);
            }
        } else {
            // Se è un numero, lo aggiungo
            risultato.push(el);
        }
    }
    return risultato;
}
\end{jsbox}

\begin{esercizio}{Es. 2: Multi-Insieme da Stringa}
\textbf{TRACCIA:} Data una stringa, creare un oggetto che rappresenti il "Multi-Insieme" dei caratteri (conteggio frequenze), ma ignorando gli spazi.

\textbf{Esempio:} `"a b a"` $\to$ `\{ "a": 2, "b": 1 \}`
\end{esercizio}

\begin{jsbox}
function contaCaratteri(testo) {
    let frequenze = {};
    
    for (let i = 0; i < testo.length; i++) {
        let char = testo[i];
        if (char === " ") continue; // Salta spazi
        
        if (frequenze[char]) {
            frequenze[char]++;
        } else {
            frequenze[char] = 1;
        }
    }
    return frequenze;
}
\end{jsbox}

\begin{esercizio}{Es. 3: Alberi Binari - Somma Foglie}
\textbf{TRACCIA:} Dato un albero binario (oggetti con `val`, `sx`, `dx`), restituire la somma dei valori delle sole \textbf{foglie} (nodi senza figli).

\textbf{Struttura:} `Nodo = { val: number, sx: Nodo | null, dx: Nodo | null }`
\end{esercizio}

\begin{jsbox}
function sommaFoglie(nodo) {
    if (!nodo) return 0;
    
    // È una foglia? (Nessun figlio)
    if (!nodo.sx && !nodo.dx) {
        return nodo.val;
    }
    
    // Altrimenti somma ricorsivamente i sottoalberi
    return sommaFoglie(nodo.sx) + sommaFoglie(nodo.dx);
}
\end{jsbox}

\begin{esercizio}{Es. 4: Intersezione di N Array}
\textbf{TRACCIA:} Scrivere una funzione che accetta un array di array (matrice) e restituisce un array contenente solo gli elementi presenti in \textbf{tutti} i sotto-array.

\textbf{Esempio:} `[[1, 2, 3], [2, 3, 4], [2, 5, 3]]` $\to$ `[2, 3]`
\end{esercizio}

\begin{jsbox}
function intersezioneTutti(matrice) {
    if (matrice.length === 0) return [];
    
    // Parto dai candidati del primo array
    let candidati = matrice[0];
    
    // Per ogni array successivo, filtro i candidati
    for (let i = 1; i < matrice.length; i++) {
        let arrayCorrente = matrice[i];
        
        // Tengo solo chi esiste anche qui
        candidati = candidati.filter(x => arrayCorrente.includes(x));
    }
    
    return candidati;
}
\end{jsbox}
 \newpage
\section{Capitolo 2: Funzioni e Programmazione Funzionale}

\begin{esercizio}{Es. 5: Pipeline di Composizione}
\textbf{TRACCIA:} Scrivere una funzione `pipeline(funzioni)` che accetta un array di funzioni $[f_1, f_2, \dots, f_n]$ e restituisce una nuova funzione che applica le $f$ in sequenza all'input.
$result = f_n(\dots f_2(f_1(x))\dots)$

\textbf{Esempio:} `[x => x+1, x => x*2]` $\to$ Input 5 $\to$ (5+1)=6 $\to$ 6*2=12.
\end{esercizio}

\begin{jsbox}
function pipeline(funzioni) {
    return function(input) {
        let valore = input;
        for (let f of funzioni) {
            valore = f(valore); // Passo il risultato alla prossima
        }
        return valore;
    };
}
// Equivalente con reduce:
// const pipeline = fs => x => fs.reduce((acc, f) => f(acc), x);
\end{jsbox}

\begin{esercizio}{Es. 6: Filtro Asimmetrico (HOF)}
\textbf{TRACCIA:} Scrivere `partition(arr, predicato)` che restituisce un oggetto `{ pass: [...], fail: [...] }`.
In `pass` vanno gli elementi per cui il predicato è true, in `fail` gli altri.

\textbf{Vincolo:} Scorrere l'array una sola volta.
\end{esercizio}

\begin{jsbox}
function partition(arr, predicato) {
    let result = { pass: [], fail: [] };
    
    for (let el of arr) {
        if (predicato(el)) {
            result.pass.push(el);
        } else {
            result.fail.push(el);
        }
    }
    return result;
}
\end{jsbox}

\begin{esercizio}{Es. 7: Generatore di Contatori (Closure)}
\textbf{TRACCIA:} Scrivere `creaContatori(n)` che restituisce un array di `n` funzioni.
La funzione all'indice `i` deve ritornare `i` quando chiamata.
\textit{Attenzione all'errore classico del `var` nei cicli!}
\end{esercizio}

\begin{jsbox}
function creaContatori(n) {
    let funzioni = [];
    
    for (let i = 0; i < n; i++) {
        // Usando 'let' nel ciclo, creiamo uno scope per ogni iterazione.
        // La closure cattura la 'i' corretta.
        funzioni.push(function() {
            return i;
        });
    }
    return funzioni;
}

// Se usassimo 'var i', tutte ritornerebbero n!
\end{jsbox}

\begin{esercizio}{Es. 8: Map su Albero (Tree Map)}
\textbf{TRACCIA:} Scrivere una funzione ricorsiva `alberoMap(nodo, f)` che restituisce un \textbf{nuovo} albero strutturalmente identico, ma dove ogni valore `val` è trasformato dalla funzione `f`.
\end{esercizio}

\begin{jsbox}
function alberoMap(nodo, f) {
    if (!nodo) return null;
    
    return {
        val: f(nodo.val),           // Trasformo il valore
        sx: alberoMap(nodo.sx, f),  // Ricorro a sinistra
        dx: alberoMap(nodo.dx, f)   // Ricorro a destra
    };
}
\end{jsbox}

\begin{esercizio}{Es. 20: Memoization (Cache Semplice)}
\textbf{TRACCIA:} Scrivere una funzione `memoize(f)` che accetta una funzione (che prende 1 argomento stringa/numero) e ne restituisce una versione "con memoria". Se richiamata con lo stesso input, non ricalcola ma restituisce il valore salvato.
\end{esercizio}

\begin{jsbox}
function memoize(f) {
    let cache = {}; // La chiusura protegge la cache
    
    return function(n) {
        if (n in cache) {
            return cache[n]; // HIT
        } else {
            let result = f(n); // MISS - Calcolo
            cache[n] = result;
            return result;
        }
    };
}
\end{jsbox}

\begin{esercizio}{Es. 21: Currying (Somma a Catena)}
\textbf{TRACCIA:} Scrivere una funzione `somma(a)` che può essere chiamata così: `somma(2)(3)(4)` e restituisce la somma finale (es. supportando 3 chiamate fisse o tornando un oggetto con `valueOf`).
\textit{Versione Semplice (3 argomenti fissi):}
\end{esercizio}

\begin{jsbox}
function somma(a) {
    return function(b) {
        return function(c) {
            return a + b + c;
        };
    };
}
// somma(1)(2)(3) -> 6
\end{jsbox}

\begin{esercizio}{Es. 22: Implementazione Manuale di 'Every'}
\textbf{TRACCIA:} Scrivere `myEvery(arr, predicato)` che torna true se \textbf{TUTTI} gli elementi soddisfano il predicato.
\textit{Appena ne trovi uno falso, devi fermarti (short-circuit).}
\end{esercizio}

\begin{jsbox}
function myEvery(arr, predicato) {
    for (let el of arr) {
        if (!predicato(el)) {
            return false; // Trovato un contromodello
        }
    }
    return true; // Nessuno ha fallito
}
\end{jsbox}

\begin{esercizio}{Es. 47: Map tramite Reduce}
\textbf{TRACCIA:} Implementare `myMap(arr, f)` usando \textbf{solo} `reduce`.
\end{esercizio}

\begin{jsbox}
function myMap(arr, f) {
    return arr.reduce((acc, el) => {
        acc.push(f(el));
        return acc; // Ricorda di ritornare l'accumulatore!
    }, []); // Valore iniziale array vuoto
}
\end{jsbox}

\begin{esercizio}{Es. 48: Filter tramite Reduce}
\textbf{TRACCIA:} Implementare `myFilter(arr, pred)` usando \textbf{solo} `reduce`.
\end{esercizio}

\begin{jsbox}
function myFilter(arr, predicato) {
    return arr.reduce((acc, el) => {
        if (predicato(el)) {
            acc.push(el);
        }
        return acc;
    }, []);
}
\end{jsbox}
 \newpage
\section{Capitolo 3: Concetti Matematici e Algoritmi}

\begin{esercizio}{Es. 9: Funzione Iniettiva (Verifica Collisioni)}
\textbf{TRACCIA:} Data una funzione `f(x)` e un dominio (array di input) `D`, scrivere `isIniettiva(f, D)` che restituisce `true` se la funzione non produce mai lo stesso output per due input diversi in D.
\end{esercizio}

\begin{jsbox}
function isIniettiva(f, dominio) {
    let immaginiViste = new Set();
    
    for (let x of dominio) {
        let y = f(x); // Calcolo immagine
        
        if (immaginiViste.has(y)) {
            return false; // Collisione trovata! (Due x danno la stessa y)
        }
        immaginiViste.add(y);
    }
    return true; // Nessuna collisione
}
\end{jsbox}

\begin{esercizio}{Es. 10: Convoluzione Discreta (Semplificata)}
\textbf{TRACCIA:} Dati due array `A` e `B` della stessa lunghezza, calcolare il prodotto scalare (somma dei prodotti elemento per elemento).
$S = \sum (A[i] \cdot B[i])$
\end{esercizio}

\begin{jsbox}
function prodottoScalare(A, B) {
    let somma = 0;
    // Assumiamo A.length === B.length
    for (let i = 0; i < A.length; i++) {
        somma += A[i] * B[i];
    }
    return somma;
}
\end{jsbox}

\begin{esercizio}{Es. 11: Punto Fisso}
\textbf{TRACCIA:} Scrivere `trovaPuntoFisso(f, min, max)` che cerca se esiste un numero intero $x$ tra `min` e `max` tale che $f(x) == x$. Se esiste ritornalo, altrimenti `null`.
\end{esercizio}

\begin{jsbox}
function trovaPuntoFisso(f, min, max) {
    for (let x = min; x <= max; x++) {
        if (f(x) === x) {
            return x; // Trovato!
        }
    }
    return null;
}
\end{jsbox}

\begin{esercizio}{Es. 12: Numeri Triangolari Ricorsivi}
\textbf{TRACCIA:} L'n-esimo numero triangolare è la somma dei numeri da 1 a n. Scrivere una funzione ricorsiva.
\end{esercizio}

\begin{jsbox}
function triangolare(n) {
    if (n <= 0) return 0;
    return n + triangolare(n - 1);
}
// Esempio: T(3) = 3 + T(2) = 3 + 2 + T(1) = 3+2+1 = 6
\begin{esercizio}{Es. 23: Numero Primo (IsPrime)}
\textbf{TRACCIA:} Scrivere `isPrime(n)` che torna true se $n$ è primo.
\textit{Ricorda: 1 non è primo. Basta controllare fino a $\sqrt{n}$.}
\end{esercizio}

\begin{jsbox}
function isPrime(n) {
    if (n < 2) return false;
    
    // Ottimizzazione: fino alla radice quadrata
    for (let i = 2; i * i <= n; i++) {
        if (n % i === 0) return false; // Divisibile
    }
    return true;
}
\end{jsbox}

\begin{esercizio}{Es. 24: MCD (Algoritmo di Euclide)}
\textbf{TRACCIA:} Calcolare il Massimo Comune Divisore tra `a` e `b`.
\textit{Algoritmo: Se b=0, MCD è a. Altrimenti MCD(b, a \% b).}
\end{esercizio}

\begin{jsbox}
function mcd(a, b) {
    if (b === 0) return a;
    return mcd(b, a % b);
}
\end{jsbox}
 \newpage
\section{Capitolo 4: TypeScript e Tipi}

\begin{esercizio}{Es. 13: Interfacce e Tipi Opzionali}
\textbf{TRACCIA:} Definire un'interfaccia `Studente` con:
- `matricola` (numero, obbligatorio)
- `nome` (stringa, obbligatorio)
- `votoLaurea` (numero, \textbf{opzionale} - potrebbe non essersi ancora laureato)

Scrivere poi una funzione `isLaureato(s)` che ritorna true se il voto esiste.
\end{esercizio}

\begin{jsbox}
interface Studente {
    matricola: number;
    nome: string;
    votoLaurea?: number; // Il punto interrogativo indica opzionale
}

function isLaureato(s: Studente): boolean {
    // Controllo se è diverso da undefined
    return s.votoLaurea !== undefined;
}
\end{jsbox}

\begin{esercizio}{Es. 14: Generics (La Scatola Magica)}
\textbf{TRACCIA:} Scrivere una classe `Scatola<T>` che può contenere un valore di qualsiasi tipo T.
Metodi: `inserisci(val: T)` e `estrai(): T`.
\end{esercizio}

\begin{jsbox}
class Scatola<T> {
    private contenuto: T | null = null;
    
    inserisci(val: T): void {
        this.contenuto = val;
    }
    
    estrai(): T | null {
        return this.contenuto;
    }
}

// Uso
let scatolaNumeri = new Scatola<number>();
scatolaNumeri.inserisci(10);
// scatolaNumeri.inserisci("ciao"); // ERRORE TS!
\end{jsbox}

\begin{esercizio}{Es. 15: Tuple e Array}
\textbf{TRACCIA:} Definire un tipo `Punto3D` come una Tripla (Tupla) di 3 numeri `[x, y, z]`.
Scrivere una funzione che somma due punti.
\end{esercizio}

\begin{jsbox}
type Punto3D = [number, number, number];

function sommaPunti(p1: Punto3D, p2: Punto3D): Punto3D {
    return [
        p1[0] + p2[0],
        p1[1] + p2[1],
        p1[2] + p2[2]
    ];
}
\end{jsbox}

\begin{esercizio}{Es. 16: Type Guards (Narrowing)}
\textbf{TRACCIA:} Una funzione accetta `input: string | number`.
Se è una stringa, ritorna la lunghezza. Se è un numero, ritorna il numero stesso moltiplicato per 2.
Usa `typeof` per distinguere.
\end{esercizio}

\begin{jsbox}
function processa(input: string | number): number {
    if (typeof input === "string") {
        // Qui TS sa che input è string
        return input.length;
    } else {
        // Qui TS sa che input e' number
        return input * 2;
    }
}
\begin{esercizio}{Es. 25: Discriminated Union (Result Type)}
\textbf{TRACCIA:} Definire un tipo `Result<T>` che può essere `{ success: true, data: T }` OPPURE `{ success: false, error: string }`.
Scrivere una funzione che gestisce entrambi i casi stampando il dato o l'errore.
\end{esercizio}

\begin{jsbox}
type Result<T> = 
    | { success: true; data: T } 
    | { success: false; error: string };

function gestisci(res: Result<number>) {
    if (res.success) {
        // Qui TS sa che esiste .data
        console.log("Dato:", res.data);
    } else {
        // Qui TS sa che esiste .error
        console.log("Errore:", res.error);
    }
}
\end{jsbox}
 \newpage
\section{Capitolo 11: La Ricorsione (Lab 11)}

\begin{dizionario}{Caso Base vs Passo Ricorsivo}
\textbf{Caso Base:} La condizione che ferma la ricorsione (es. `n == 0`). Senza di esso = Stack Overflow.
\textbf{Passo Ricorsivo:} La funzione chiama se stessa con un problema *più piccolo*.
\end{dizionario}

\begin{esercizio}{Es. 26: Somma Numeri in Array (Ricorsiva)}
\textbf{TRACCIA:} Scrivere `sumRecursive(arr)` che somma gli elementi di un array usando la ricorsione (senza cicli for/while).
\textit{Suggerimento: somma = testa + somma(coda).}
\end{esercizio}

\begin{jsbox}
function sumRecursive(arr) {
    // Caso Base: Array vuoto -> somma 0
    if (arr.length === 0) return 0;
    
    // Passo Ricorsivo: Primo elemento + somma del resto
    return arr[0] + sumRecursive(arr.slice(1));
}
\end{jsbox}

\begin{esercizio}{Es. 27: Palindroma Ricorsiva}
\textbf{TRACCIA:} Scrivere `isPalindrome(str)` che controlla se una stringa è palindroma ricorsivamente.
\textit{Logica: Prima lettera == Ultima lettera AND il centro è palindromo.}
\end{esercizio}

\begin{jsbox}
function isPalindrome(str) {
    // Casi Base
    if (str.length <= 1) return true;
    
    // Controllo estremi
    if (str[0] !== str[str.length - 1]) return false;
    
    // Ricorsione sul centro
    return isPalindrome(str.slice(1, -1));
}
\end{jsbox}

\begin{esercizio}{Es. 28: Tartaglia (Coefficiente Binomiale)}
\textbf{TRACCIA:} Calcolare $C(n, k)$ (Tartaglia) ricorsivamente.
Formula: $C(n, k) = C(n-1, k-1) + C(n-1, k)$.
\end{esercizio}

\begin{jsbox}
function tartaglia(n, k) {
    if (k === 0 || k === n) return 1; // Bordi del triangolo
    if (k > n) return 0;
    
    return tartaglia(n - 1, k - 1) + tartaglia(n - 1, k);
}
\end{jsbox}
 \newpage
\section{Capitoli 12-14: Liste Concatenate}

\begin{dizionario}{Lista vs Array}
\textbf{Lista:} Nodi sparsi collegati da punti `next`. Accesso $O(N)$.
\textbf{Array:} Blocco contiguo. Accesso $O(1)$.
\textbf{Nodo:} Oggetto `{ val: ..., next: ... }`.
\end{dizionario}

\begin{esercizio}{Es. 29: Conta Occorrenze (Lista)}
\textbf{TRACCIA:} Scrivere `countVal(nodo, target)` che conta quante volte `target` appare nella lista.
\end{esercizio}

\begin{jsbox}
function countVal(nodo, target) {
    if (!nodo) return 0;
    
    let match = (nodo.val === target) ? 1 : 0;
    return match + countVal(nodo.next, target);
}
\end{jsbox}

\begin{esercizio}{Es. 30: Ultimo Elemento (Il 'Last')}
\textbf{TRACCIA:} Restituire il valore dell'ultimo nodo della lista.
\textit{Attenzione: non l'intera lista, solo il valore.}
\end{esercizio}

\begin{jsbox}
function getLast(nodo) {
    if (!nodo) return null; // Lista vuota
    if (!nodo.next) return nodo.val; // Sono l'ultimo!
    return getLast(nodo.next);
}
\end{jsbox}

\begin{esercizio}{Es. 31: Filtra Lista (List Filter)}
\textbf{TRACCIA:} Scrivere `listFilter(nodo, predicato)` che restituisce una **nuova** lista contenente solo i nodi che passano il test.
\end{esercizio}

\begin{jsbox}
function listFilter(nodo, predicato) {
    if (!nodo) return null;
    
    // Se il nodo corrente passa, lo tengo e continuo
    if (predicato(nodo.val)) {
        return {
            val: nodo.val,
            next: listFilter(nodo.next, predicato)
        };
    } else {
        // Se non passa, lo salto (ritorno direttamente il risultato del prossimo)
        return listFilter(nodo.next, predicato);
    }
}
\end{jsbox}

\begin{esercizio}{Es. 32: Append (Concatena Due Liste)}
\textbf{TRACCIA:} Scrivere `listConcat(l1, l2)` che attacca `l2` in fondo a `l1`.
\end{esercizio}

\begin{jsbox}
function listConcat(l1, l2) {
    if (!l1) return l2; // Se l1 è finita, attacco l2
    
    return {
        val: l1.val,
        next: listConcat(l1.next, l2)
    };
}
\end{jsbox}

\begin{esercizio}{Es. 42: Inverti Lista (Reverse)}
\textbf{TRACCIA:} Scrivere `reverseList(head)` che inverte i puntatori della lista e restituisce la nuova testa.
\textit{Iterativo o Ricorsivo.}
\end{esercizio}

\begin{jsbox}
function reverseList(head) {
    let prev = null;
    let curr = head;
    
    while (curr !== null) {
        let nextTemp = curr.next; // Salvo il prossimo
        curr.next = prev;         // Inverto il puntatore
        prev = curr;              // Avanzo prev
        curr = nextTemp;          // Avanzo curr
    }
    return prev; // Nuova testa
}
\end{jsbox}

\begin{esercizio}{Es. 43: Rileva Ciclo (Floyd)}
\textbf{TRACCIA:} Determinare se una lista ha un ciclo (un nodo punta a uno precedente).
\textit{Usa due puntatori: 'tartaruga' (lento) e 'lepre' (veloce).}
\end{esercizio}

\begin{jsbox}
function hasCycle(head) {
    let slow = head;
    let fast = head;
    
    while (fast && fast.next) {
        slow = slow.next;       // 1 passo
        fast = fast.next.next;  // 2 passi
        
        if (slow === fast) return true; // Si sono incontrati!
    }
    return false;
}
\end{jsbox}
 \newpage
\section{Capitoli 16-17: Array Avanzati e Destructuring}

\begin{dizionario}{Destructuring}
\textbf{Sintassi:} `const {a, b} = obj` oppure `const [x, y] = arr`.
Permette di estrarre valori in modo conciso.
\end{dizionario}

\begin{esercizio}{Es. 33: Swap Variabili (Destructuring)}
\textbf{TRACCIA:} Scambiare i valori di due variabili `a` e `b` senza usare una variabile temporanea `temp`, usando il destructuring.
\end{esercizio}

\begin{jsbox}
let a = 1, b = 2;
[a, b] = [b, a];
console.log(a, b); // 2, 1
\end{jsbox}

\begin{esercizio}{Es. 34: Somma Variadica (Rest Operator)}
\textbf{TRACCIA:} Scrivere una funzione `sommaTutto(...numeri)` che accetta un numero arbitrario di argomenti e li somma tutti.
\end{esercizio}

\begin{jsbox}
function sommaTutto(...numeri) {
    // 'numeri' è un vero array [n1, n2, ...]
    return numeri.reduce((acc, n) => acc + n, 0);
}
console.log(sommaTutto(1, 2, 3, 4)); // 10
\end{jsbox}

\begin{esercizio}{Es. 35: Merge di Oggetti (Spread)}
\textbf{TRACCIA:} Dati due oggetti `defaultConfig` e `userConfig`, creare un `finalConfig` che ha tutte le proprietà, ma `userConfig` vince in caso di conflitto.
\end{esercizio}

\begin{jsbox}
const defaultConfig = { theme: "light", debug: false };
const userConfig = { theme: "dark" };

const finalConfig = { ...defaultConfig, ...userConfig };
// { theme: "dark", debug: false }
\end{jsbox}

\begin{esercizio}{Es. 46: Rotazione Matrice 90 Gradi}
\textbf{TRACCIA:} Ruotare una matrice $N \times N$ di 90 gradi in senso orario, \textbf{in-place} (modificando la matrice originale).
\end{esercizio}

\begin{jsbox}
function rotate(matrix) {
    let n = matrix.length;
    
    // 1. Trasposta
    for (let i = 0; i < n; i++) {
        for (let j = i; j < n; j++) {
            [matrix[i][j], matrix[j][i]] = [matrix[j][i], matrix[i][j]];
        }
    }
    
    // 2. Inverti ogni riga
    for (let i = 0; i < n; i++) {
        matrix[i].reverse();
    }
}
\end{jsbox}
 \newpage
\section{Capitolo 18: Alberi Binari e BST}

\begin{dizionario}{BST (Binary Search Tree)}
Albero dove per ogni nodo:
- Tutti i valori a \textbf{Sinistra} sono minori.
- Tutti i valori a \textbf{Destra} sono maggiori.
Permette ricerche veloci (come la ricerca binaria).
\end{dizionario}

\begin{esercizio}{Es. 36: Conteggio Nodi}
\textbf{TRACCIA:} Scrivere una funzione ricorsiva che conta quanti nodi ci sono in totale nell'albero.
\end{esercizio}

\begin{jsbox}
function contaNodi(root) {
    if (!root) return 0;
    return 1 + contaNodi(root.sx) + contaNodi(root.dx);
}
\end{jsbox}

\begin{esercizio}{Es. 37: Altezza Albero}
\textbf{TRACCIA:} Calcolare l'altezza dell'albero (il percorso più lungo dalla radice a una foglia).
\end{esercizio}

\begin{jsbox}
function altezza(root) {
    if (!root) return -1; // O 0, dipende dalla definizione. Qui -1 per albero vuoto.
    
    let hSx = altezza(root.sx);
    let hDx = altezza(root.dx);
    
    return 1 + Math.max(hSx, hDx);
}
\end{jsbox}

\begin{esercizio}{Es. 38: Verifica BST}
\textbf{TRACCIA:} Controllare se un albero è un BST valido (ogni nodo deve rispettare min/max ereditati dai genitori).
\end{esercizio}

\begin{jsbox}
function isBST(nodo, min = -Infinity, max = Infinity) {
    if (!nodo) return true;
    
    if (nodo.val <= min || nodo.val >= max) return false;
    
    return isBST(nodo.sx, min, nodo.val) && 
           isBST(nodo.dx, nodo.val, max);
}
\begin{esercizio}{Es. 44: Specchia Albero (Mirror)}
\textbf{TRACCIA:} Scrivere una funzione che inverte speculare (sx diventa dx) tutto l'albero.
\end{esercizio}

\begin{jsbox}
function invertTree(root) {
    if (!root) return null;
    
    // Scambio i figli
    let temp = root.sx;
    root.sx = root.dx;
    root.dx = temp;
    
    // Ricorsione
    invertTree(root.sx);
    invertTree(root.dx);
    
    return root;
}
\end{jsbox}

\begin{esercizio}{Es. 45: Albero Simmetrico}
\textbf{TRACCIA:} Verificare se un albero è speculare a se stesso (il sottoalbero sx è il mirror del dx).
\end{esercizio}

\begin{jsbox}
function isSymmetric(root) {
    if (!root) return true;
    
    function isMirror(n1, n2) {
        if (!n1 && !n2) return true; // Entrambi null
        if (!n1 || !n2) return false; // Uno solo null
        
        return (n1.val === n2.val) &&
               isMirror(n1.sx, n2.dx) &&
               isMirror(n1.dx, n2.sx);
    }
    
    return isMirror(root.sx, root.dx);
}
\end{jsbox}
 \newpage
\section{Capitoli 20-21: Grafi}

\begin{dizionario}{Grafo Orientato vs Non Orientato}
\textbf{Orientato:} A $\to$ B (Senso unico).
\textbf{Non Orientato:} A $-$ B (Doppio senso).
\textbf{Ciclo:} Un percorso che parte da un nodo e torna allo stesso nodo.
\end{dizionario}

\begin{esercizio}{Es. 39: Grado di un Nodo}
\textbf{TRACCIA:} Dato un grafo (mappa di adiacenza) e un nodo, restituire il numero di vicini (grado).
\end{esercizio}

\begin{jsbox}
function grado(grafo, nodo) {
    if (!grafo.has(nodo)) return 0;
    return grafo.get(nodo).length;
}
\end{jsbox}

\begin{esercizio}{Es. 40: Esiste Percorso (DFS)}
\textbf{TRACCIA:} Scrivere `haPercorso(g, start, end)` che ritorna true se esiste una strada da `start` a `end`.
Usa DFS e un Set `visited` per evitare loop.
\end{esercizio}

\begin{jsbox}
function haPercorso(grafo, curr, target, visited = new Set()) {
    if (curr === target) return true;
    
    visited.add(curr);
    let vicini = grafo.get(curr) || [];
    
    for (let v of vicini) {
        if (!visited.has(v)) {
            if (haPercorso(grafo, v, target, visited)) return true;
        }
    }
    return false;
}
\end{jsbox}

\begin{esercizio}{Es. 41: Distanza Minima (BFS)}
\textbf{TRACCIA:} Trovare la distanza minima (numero di passi) tra `start` e `end`.
Usa BFS (Code).
\end{esercizio}

\begin{jsbox}
function distanzaMinima(grafo, start, end) {
    let queue = [[start, 0]]; // [nodo, distanza]
    let visited = new Set([start]);
    
    while (queue.length > 0) {
        let [curr, dist] = queue.shift();
        
        if (curr === end) return dist;
        
        let vicini = grafo.get(curr) || [];
        for (let v of vicini) {
            if (!visited.has(v)) {
                visited.add(v);
                queue.push([v, dist + 1]);
            }
        }
    }
    return -1; // irraggiungibile
}
\end{jsbox}





\section{L'identità segreta degli oggetti}

\subsection{Introduzione e Contenuti}
Il laboratorio si focalizza sulla comprensione profonda degli oggetti in JavaScript, analizzando la differenza tra valore e riferimento, l'uso di costruttori, metodi, prototipi e classi.

\subsection{Valori e Riferimenti}
JavaScript gestisce i dati in due modi fondamentali in base alla loro natura (base o complessa).

\begin{definition}[Tipi di dato e Semantica]
    \begin{itemize}
        \item \textbf{Tipi Base:} (numeri, booleani, stringhe). Hanno una \textbf{semantica per valore}. Le variabili contengono direttamente il dato.
        \item \textbf{Tipi Complessi:} (array, oggetti, funzioni). Hanno una \textbf{semantica per riferimento}. Le variabili contengono un indirizzo di memoria (riferimento).
    \end{itemize}
\end{definition}

\subsection{Meccanismi di Assegnamento e Uguaglianza}
\begin{itemize}
    \item \textbf{Per Valore:} L'assegnamento copia il valore; l'uguaglianza confronta il contenuto.
    \item \textbf{Per Riferimento:} L'assegnamento copia il riferimento; l'uguaglianza (\texttt{==}) confronta se i riferimenti puntano alla stessa area di memoria (identità), non se i contenuti sono uguali.
\end{itemize}

\begin{jsbox}[Esempio Identità vs Uguaglianza]
var a = [1,2,3]; // Riferimento alfa
var b = [1,2,3]; // Riferimento beta
var c = a;       // Copia riferimento alfa

a == b; // false (alfa != beta)
a == c; // true  (alfa == alfa)
\end{jsbox}

\begin{concept}[Allocazione di Memoria]
Ogni volta che viene valutato un letterale (es. \texttt{[1,2,3]} o \texttt{\{...\}}), viene allocata una \textbf{nuova area di memoria} con un nuovo riferimento unico.
\end{concept}

\section{Modifica di Riferimenti e Funzioni}
La condivisione di oggetti in memoria implica che modifiche effettuate tramite un riferimento siano visibili attraverso tutti gli altri puntatori allo stesso oggetto.

\subsection{Passaggio di argomenti}
JavaScript usa il \textbf{passaggio per valore} degli argomenti. Tuttavia:
\begin{itemize}
    \item Se passiamo un tipo base, la funzione lavora su una copia locale e non può modificare la variabile originale.
    \item Se passiamo un riferimento, la funzione riceve una copia del riferimento. Può quindi modificare le proprietà dell'oggetto originale.
\end{itemize}

\begin{jsbox}[Modifica Oggetto Riferito]
function compleanno(p) { p.età++; }
var pippo = {nome: "Pippo", età: 35}; // Riferimento alfa
compleanno(pippo); // Passa alfa
// pippo -> {nome: "Pippo", età: 36}
\end{jsbox}

\section{Strutture Dati Complesse: Grafi}
A differenza degli alberi, dove ogni nodo compare in un solo posto (proprietà di unicità del padre), nei \textbf{grafi} lo stesso oggetto può comparire in più campi.

\begin{definition}[Identità nel Grafo]
In un grafo $G = \langle N, E \rangle$:
\begin{itemize}
    \item \textbf{Nodi:} oggetti del tipo \texttt{\{ val: x \}}.
    \item \textbf{Archi:} oggetti \texttt{\{ da: n1, a: n2 \}} dove \texttt{n1} e \texttt{n2} sono riferimenti ai nodi.
\end{itemize}
Per far sì che più archi puntino allo \textbf{stesso} nodo (identità), si usano i riferimenti alle variabili.
\end{definition}

\section{Funzioni, Metodi e Costruttori}

\subsection{Sintassi abbreviata dei Metodi}
Esistono tre modi per definire funzioni come proprietà di un oggetto:
\begin{jsbox}[Sintassi Metodi]
var o = {
    inc: x => x + 1,                    // Arrow function
    dec: function(x) { return x - 1; }, // Classica
    zero(x) { return x - x; }           // Abbreviata (Metodo)
};
\end{jsbox}

\subsection{Oggetti Globali Pre-dichiarati}
Variabili come \texttt{console} e \texttt{Math} sono \textbf{oggetti globali}.
\begin{itemize}
    \item \texttt{typeof console} $\rightarrow$ \texttt{"object"}.
    \item \texttt{Math.max()} o \texttt{console.log()} sono proprietà di tipo funzione di questi oggetti.
\end{itemize}

\begin{curiosita}[Dinamicità]
È possibile aggiungere funzioni personalizzate a questi oggetti (es. \texttt{Math.fattoriale = ...}) per usarle nel resto del programma.
\end{curiosita}

\subsection{La parola chiave \texttt{this}}
Quando una funzione è proprietà di un oggetto e vi accede internamente, viene chiamata \textbf{metodo}. Si usa \texttt{this} per riferirsi all'oggetto corrente.

\begin{jsbox}[Uso di this]
persona.compleanno = function() { this.età++; };
pippo.compleanno(); // 'this' punta a pippo
\end{jsbox}

\subsection{L'operatore \texttt{new}}
Crea un nuovo oggetto vuoto, imposta \texttt{this} a tale oggetto nel costruttore e lo restituisce implicitamente.
\begin{concept}[Convenzione Costruttori]
Le funzioni destinate a \texttt{new} iniziano con la \textbf{lettera maiuscola}.
\end{concept}

\section{Il Mistero dei Prototipi}
Perché un array ha il metodo \texttt{.push()} se non è definito direttamente nell'oggetto?

\begin{definition}[Meccanismo di Ricerca]
Ogni oggetto ha un \textbf{prototipo}. 
\begin{enumerate}
    \item Se cerco una chiave e non c'è nell'oggetto, JavaScript la cerca nel prototipo.
    \item Si risale la \textbf{catena dei prototipi} fino a \texttt{Object.prototype}.
    \item Se la chiave non esiste in tutta la catena, il risultato è \texttt{undefined}.
\end{enumerate}
\end{definition}

Il prototipo di un oggetto si ottiene con \texttt{Object.getPrototypeOf(o)} o tramite la proprietà speciale \texttt{\_\_proto\_\_}. L'oggetto prototipo degli array contiene tutti i metodi standard: \textit{push, pop, map, filter, etc.}.

\section{Conclusione: Verso le Classi}
Il meccanismo dei prototipi e dei costruttori è la base per l'uso delle \textbf{Classi} in JavaScript, tema centrale della prossima lezione.
 \newpage
\section{L'identità segreta degli oggetti: Prototipi e Classi}

\subsection{Recap: Valori, Riferimenti e Indirezione}
In JavaScript, la distinzione tra tipi base e complessi è fondamentale:
\begin{itemize}
    \item \textbf{Tipi Base:} (booleani, numeri, stringhe) $\rightarrow$ Semantica per \textbf{valore}.
    \item \textbf{Tipi Complessi:} (oggetti, array, funzioni) $\rightarrow$ Semantica per \textbf{riferimento}.
\end{itemize}

La semantica per riferimento utilizza indirizzi di memoria come valori. L'accesso ai dati avviene tramite operazioni di \textbf{indirezione} (spesso implicite in JS).

\begin{definition}[Operazioni di Indirezione]
    \begin{itemize}
        \item \texttt{oggetto.chiave}: Va in memoria all'indirizzo dell'oggetto, cerca la chiave e restituisce il contenuto.
        \item \texttt{array[indice]}: Analogo al precedente, usando l'indice come chiave.
        \item \texttt{funzione(...)}: Va in memoria a recuperare la definizione ed esegue il codice.
    \end{itemize}
\end{definition}

\section{Identità vs Uguaglianza: Il caso dei Grafi}
L'uso dei riferimenti è cruciale nelle strutture dati complesse come i grafi, dove lo \textbf{stesso} oggetto nodo può essere puntato da molteplici archi.

\begin{concept}[Identità]
Dire "lo stesso oggetto" (Identità) non è come dire "un oggetto uguale" (Uguaglianza).
\begin{itemize}
    \item \textbf{Uguaglianza:} Due oggetti distinti con lo stesso contenuto (es. \texttt{\{val:1\}} e \texttt{\{val:1\}}).
    \item \textbf{Identità:} Due riferimenti che puntano alla stessa area di memoria.
\end{itemize}
\end{concept}

Per modellare un grafo $G=\langle N,E\rangle$ dove più archi insistono sullo stesso nodo, è indispensabile usare variabili che contengano i riferimenti ai nodi, e riutilizzare quelle variabili nella creazione degli archi.

\section{Costruttori e l'Operatore \texttt{new}}
Per garantire che gli oggetti abbiano sempre la struttura corretta (campi giusti), si usano funzioni \textbf{costruttori}.

\subsection{Funzionamento di \texttt{new}}
L'operatore \texttt{new}, anteposto a una chiamata di funzione, altera il comportamento standard:
\begin{enumerate}
    \item \textbf{Allocazione:} Crea un nuovo oggetto vuoto (chiamiamolo \textit{o}).
    \item \textbf{Invocazione:} Chiama la funzione costruttore impostando \texttt{this = o}.
    \item \textbf{Arricchimento:} Il codice della funzione usa \texttt{this} per aggiungere proprietà e metodi a \textit{o}.
    \item \textbf{Ritorno:} Restituisce implicitamente \textit{o}.
\end{enumerate}

\begin{jsbox}[Confronto Costruttori]
// 1. Factory Function (Sconsigliato - crea oggetti, non istanze legate)
function grafo(N, E) {
    return { nodi: N, archi: E, size() { return this.nodi.length } };
}

// 2. Constructor Function con 'new' (Standard pre-ES6)
function Grafo(N, E) {
    this.nodi = N;
    this.archi = E;
    // this è l'oggetto appena creato da new
}
var g = new Grafo([], []);
\end{jsbox}

\section{I Prototipi}
Il meccanismo dei prototipi permette agli oggetti di ereditare funzionalità senza duplicarle in memoria.

\begin{definition}[Algoritmo di Lookup delle Proprietà]
Quando si accede a \texttt{obj.key}:
\begin{enumerate}
    \item JS cerca \texttt{key} direttamente in \texttt{obj}. Se trovata, restituisce il valore.
    \item Se non trovata, cerca nel \textbf{prototipo} di \texttt{obj}.
    \item Continua a risalire la catena dei prototipi.
    \item Se arriva a un prototipo \texttt{null} (fine catena) e non ha trovato nulla, restituisce \texttt{undefined}.
\end{enumerate}
\textit{Nota: La scrittura avviene sempre sull'oggetto stesso, mai sul prototipo.}
\end{definition}

\subsection{Visualizzazione della Catena}
Possiamo ispezionare il prototipo con \texttt{Object.getPrototypeOf(o)} o la proprietà deprecata \texttt{\_\_proto\_\_}.

\begin{center}
\begin{tikzpicture}[
    object/.style={rectangle, draw=paletteIndigo, fill=paletteLavender!20, thick, minimum width=2.5cm, minimum height=1.5cm, align=center},
    pointer/.style={->, >=stealth, thick, color=paletteRegalia},
    node distance=1.5cm and 2cm
]

% Oggetto array
\node[object] (arr) {Array \texttt{a} \\ \texttt{[1, 5, 7]} \\ \footnotesize \texttt{\_\_proto\_\_: ?}};

% Array Prototype
\node[object, right=of arr] (arrProto) {\textbf{Array.prototype} \\ \textit{push, pop, map...} \\ \footnotesize \texttt{\_\_proto\_\_: ?}};

% Object Prototype
\node[object, right=of arrProto] (objProto) {\textbf{Object.prototype} \\ \textit{toString, valueOf...} \\ \footnotesize \texttt{\_\_proto\_\_: null}};

% Null
\node[draw=none, right=of objProto] (null) {\textbf{null}};

% Collegamenti
\draw[pointer] (arr.east) -- node[above]{\footnotesize prototipo} (arrProto.west);
\draw[pointer] (arrProto.east) -- node[above]{\footnotesize prototipo} (objProto.west);
\draw[pointer] (objProto.east) -- (null.west);

\end{tikzpicture}
\end{center}

\begin{warning}[Enumerabilità]
Il ciclo \texttt{for (k in o)} e \texttt{Object.keys(o)} restituiscono solo le chiavi \textbf{proprie} dell'oggetto, ignorando quelle ereditate dal prototipo.
\end{warning}

\section{Prototipi e Costruttori: Il Doppio Salto Carpiato}
Esiste un legame simbiotico tra una funzione costruttore e i prototipi degli oggetti che essa crea.
\begin{concept}[Proprietà \texttt{.prototype}]
Ogni funzione \texttt{F} ha una proprietà \texttt{F.prototype}. Quando chiamiamo \texttt{new F()}, l'oggetto creato riceve automaticamente \texttt{F.prototype} come suo prototipo interno (\texttt{\_\_proto\_\_}).
\end{concept}

Questo permette di definire metodi una sola volta nel prototipo del costruttore, rendendoli disponibili a tutte le istanze:
\begin{jsbox}[Pattern Costruttore + Prototipo]
function Persona(n, e) {
    this.nome = n; // Stato interno (diverso per ogni istanza)
    this.età = e;
}
// Comportamento condiviso (uguale per tutte le istanze)
Persona.prototype.compleanno = function() { this.età++; };

var p = new Persona("Pippo", 35);
// p.__proto__ === Persona.prototype è TRUE
\end{jsbox}

\section{Le Classi (Syntactic Sugar)}
Le versioni recenti di JavaScript introducono la keyword \texttt{class}. Non è un nuovo modello a oggetti, ma una sintassi più pulita sopra il sistema prototipale esistente.

\begin{jsbox}[Sintassi Class]
class Persona {
    constructor(nome, eta) {
        this.nome = nome;
        this.età = eta;
    }

    compleanno() {
        this.età++;
    }
}
\end{jsbox}

\subsection{Cosa accade dietro le quinte?}
Anche usando \texttt{class}:
\begin{itemize}
    \item \texttt{typeof Persona} restituisce \textbf{"function"}.
    \item I metodi definiti nella classe (es. \texttt{compleanno}) finiscono in \texttt{Persona.prototype}.
    \item Le proprietà definite nel costruttore (es. \texttt{nome}) finiscono direttamente nell'istanza oggetto.
\end{itemize}

\section{La "Mappa Proibita" (Struttura Profonda di JS)}
Le slide finali mostrano la struttura ricorsiva e intrecciata di JS. Ecco la logica:

\begin{itemize}
    \item \textbf{Ogni funzione è un oggetto:} Quindi \texttt{Array} (il costruttore) ha come prototipo \texttt{Function.prototype}.
    \item \textbf{Function.prototype è un oggetto:} Quindi il suo prototipo è \texttt{Object.prototype}.
    \item \textbf{Object è una funzione:} Il suo prototipo è \texttt{Function.prototype}.
\end{itemize}

\begin{center}
\begin{tikzpicture}[
    node distance=2cm,
    entity/.style={rectangle, draw=paletteRegalia, fill=white, thick, minimum size=1cm},
    proto/.style={->, >=stealth, dashed, color=gray!80}
]

\node[entity] (funcProto) {Function.prototype};
\node[entity, below=of funcProto] (objProto) {Object.prototype};
\node[entity, right=of funcProto] (ArrayConst) {Array (Funzione)};
\node[entity, below=of ArrayConst] (ArrProto) {Array.prototype};
\node (null) [right=of objProto] {null};

% Relazioni
\draw[proto] (ArrayConst) -- node[right]{\tiny \_\_proto\_\_} (funcProto);
\draw[proto] (funcProto) -- node[right]{\tiny \_\_proto\_\_} (objProto);
\draw[proto] (ArrProto) -- node[right]{\tiny \_\_proto\_\_} (objProto);
\draw[proto] (objProto) -- (null);

% Relazione prototype
\draw[->, thick, color=paletteIndigo] (ArrayConst) -- node[left]{\tiny .prototype} (ArrProto);

\end{tikzpicture}
\end{center}


\end{document}
